\subsection{Countable internal gates and finite-diagram replacements}

The familiar countable-subalgebra construction appears in the proof
of the Stacks Project, Lemma~22.38.3, tag~0CRM
\cite{Root_CountableDGA}. We opened that proof. Here we prove the
additional assertions required for both gradings, divided powers,
the prescribed first skeleton, every eventual exactness range, and
finite diagrams of actual morphisms. We claim no priority for the
countable-closure method.

An object is called countable here when the disjoint union of all
its groups $A_n^{w,q}$, as underlying sets, is countable. Since
$(n,w,q)$ ranges over $\mathbb N^3$, this is equivalent to
countability of each group. Write
\[
 \mathcal C^{[d]}_\omega\subset\mathcal C^{[d]},\qquad
 \mathcal E^{[d]}_\omega\subset\mathcal E^{[d]}
 \tag{S09.155}
\]
for the full subcategories with this property. Coefficients remain
$S=\Gamma_{\mathbb Z}(T)$; no scalar extension or completion is used.

\begin{theorem}[A-S09-29: Proved; countable internal reductions]
\label{AS0929}
For every fixed $d\geq0$, each inclusion in (S09.155) induces
an equivalence on all-arrow localizations. Both inclusions are
proper, even up to strict isomorphism.

The common-source assertion for $\mathcal A_{\mathbb Z}$ is
equivalent to common sources with targets and source internal to
$\mathcal E^{[d]}_\omega$. In its restricted test on target pairs,
their first skeletons may be the same $P(N)$, while the source
has its own cutoff.

The terminal-localization assertion for $\mathcal A_{\mathbb Z}$
is equivalent to terminality of the all-arrow localization of
$\mathcal C^{[d]}_\omega$. Every statement holds with support
$q\leq w$, termwise flatness over $\mathbb Z$, or both.
The countable common-source and terminal-localization questions
themselves remain unresolved.
\end{theorem}

\begin{proof}
We prove a finite-diagram assertion that supplies all the required
maps. Let $\mathcal D$ be either $\mathcal C^{[d]}$ or
$\mathcal E^{[d]}$. Take finitely many actual objects $A_i$ in
$\mathcal D$ and finitely many specified actual arrows between
them. All equations between composites of these arrows are
retained as equations in the ambient category. At each vertex
specify a countable set of homogeneous elements to be included.
Some vertices may themselves be countable and required to remain
unchanged; include every element of those vertices in the specified
sets. We will construct actual subobjects and restricted arrows
\[
 A'_i\hookrightarrow A_i,\qquad
 f|_{A'_i}:A'_i\longrightarrow A'_j\quad(f:A_i\to A_j),
 \qquad A'_i\in\mathcal D_\omega.
 \tag{S09.156}
\]
The inclusions will commute strictly with every specified arrow.

Choose nondecreasing admissibility bounds $\lambda_i(q),\mu_i(q)$
for each original object, by replacing any given bounds with the
maximum of their values up to $q$. The coefficient comparison
through $d$ already holds in all weights. In the
$\mathcal E^{[d]}$ case also retain its specified cutoff $N_i$
for rows through $d+2$ and its first-skeleton presentation $P(N_i)$.
In the $\mathcal C^{[d]}$ case choose a strict first-skeleton
identification with $P(E_i)$ from A-006, and an integer $N_i$
such that every set $(E_i)_w$ is nonempty for $w\geq N_i$.
These choices are made for this finite diagram. No functorial
choice on the entire category is required.

We first describe a countable closure operation on the finite
collection of subsets. Adjoin all coefficient elements in every
simplex. Close under addition, additive inverses, multiplication,
the vertical differential, all defined divided powers, every face
and degeneracy, and every specified arrow. There are countably
many finitary operations and countably many indices. Applying
them in countably many stages to countable sets gives countable
sets: at each stage a countable union of finite Cartesian powers
of countable sets is countable. The resulting sets form strict
simplicial PD differential subalgebras and a subdiagram. Weight
and form components are retained separately throughout. The
augmentation is the restriction of the original augmentation.

We impose the first-skeleton condition during this same closure.
In the $\mathcal E^{[d]}$ case include the whole of $(A_i)_0$
and $(A_i)_1$ from the beginning. This is countable, since
$P(N_i)$ is a free PD algebra on countably many generators over
the countable ring $S$. Its inclusion makes the resulting first
skeleton literally equal to that of $A_i$.

For $\mathcal C^{[d]}$, initially choose one index in $(E_i)_w$
for each $w\geq N_i$. Whenever an element of simplex degree one
is added, write its inverse image under the chosen presentation
as its finite free PD monomial expansion. Include every variable
index occurring in that expansion, together with its differential
partner. Include the entire first-skeleton subalgebra on these
indices. This is another countable closure operation: each
polynomial has finitely many indices, and the free PD algebra on
countably many indices is countable. Iterating this rule together
with all the preceding rules produces in simplex one exactly
the free subalgebra on the selected union $F_i\subset E_i$.
Its intersection with simplex zero is $S$ and its faces are the
specified coefficient endpoints. Hence its first skeleton is
strictly $P(F_i)$, with $(F_i)_w$ nonempty for all $w\geq N_i$.
The independence of its monomials follows from their inclusion
among the monomial basis in $P(E_i)$; no new algebra relation
is imposed. Countable vertices that were fixed initially remain
fixed during this construction.

We now add the exactness witnesses. Start with the closure just
described, and denote its vertex subalgebras by $B_i^0$. At stage
$r$, perform both of the following choices for every vertex.

For every homogeneous vertically closed element
$z\in(B_i^r)_n^{w,q}$ with $q>0$, require a primitive when
\[
                 q\leq d\quad\hbox{or}\quad w\geq\lambda_i(q).
 \tag{S09.157}
\]
The same element is closed in $A_i$. Its cohomology class there
is zero in exactly the required ranges, so choose
$v_z\in(A_i)_n^{w,q-1}$ with $dv_z=z$ and adjoin it.

For every homogeneous normalized horizontal cycle
$z\in N_n((B_i^r)^{w,q})$, including all elements for $n=0$,
require a normalized primitive when
\[
 w\geq\mu_i(q),\quad\hbox{or, in the $\mathcal E^{[d]}$ case,}
 \quad q\leq d+2\ \hbox{and}\ w\geq N_i.
 \tag{S09.158}
\]
The cycle is an actual normalized cycle in $A_i$ and the
indicated row there is exact. Choose
$h_z\in N_{n+1}(A_i^{w,q})$ with $d_0h_z=z$ and adjoin it.
The equations $d_jh_z=0$ for $j>0$ are part of this choice.
These are choices of elements of the original algebra, with
their actual weights and form degrees.

There are countably many choices at each stage, since all
indices and all elements in $B_i^r$ form a countable set.
Apply the full algebra, diagram and first-skeleton closure to
these added witnesses and call the result $B_i^{r+1}$.
Define
\[
                         A'_i=\bigcup_{r\geq0} B_i^r.
 \tag{S09.159}
\]
These are actual countable subalgebras with the strict maps
(S09.156). Every polynomial identity and every equality of
composites is inherited from the ambient algebra and diagram.
All operations are finitary, so their values belong to the
union. No completed algebra or infinite sum has been introduced.

We check admissibility directly. A closed positive form in one
of the ranges (S09.157) belongs to one stage and its selected
primitive belongs to the next. Thus its cohomology class in
$A'_i$ vanishes. For zero-forms, the original global degree-zero
coefficient identification in $A_i$ says that every closed
zero-form is the image of one and only one element of $S$.
All coefficients were included in $A'_i$ and the inclusion is
injective, so the same identification holds in $A'_i$ in every
weight and simplex. This checks its comparison through $d$
globally and its eventual vertical comparison in all degrees.
Nondecreasing $\lambda_i$ ensures that the stated bound for a
degree $e$ controls all $q\leq e$.

Likewise every horizontal normalized cycle in (S09.158) lies
in one stage, and its normalized primitive lies in the next.
Normalization here means the intersection of the finitely many
kernels $d_j$, $1\leq j\leq n$, inside the given simplex group;
the equations are unchanged under our inclusions. Thus this
argument includes degree-zero homology and all simplex degrees
with the original uniform cutoffs. It proves eventual row
exactness, and in the $\mathcal E^{[d]}$ case the entire row
bound $N_i$ through $d+2$. Together with the preserved first
skeleton this proves $A'_i\in\mathcal D_\omega$.

The support property is inherited by subgroups. If all original
terms are flat over $\mathbb Z$, they are torsionfree; their
subgroups are torsionfree and hence flat over $\mathbb Z$.
This proves the finite-diagram assertion in every stated variant.

We next prove the localization statement. Every object $A$ of
$\mathcal D$ has an actual inclusion $A'\hookrightarrow A$
with $A'\in\mathcal D_\omega$, by the one-vertex case of
(S09.156). Thus the inclusion of localizations is essentially
surjective. A word of arrows and formal inverses between two
countable objects uses a finite diagram. Apply (S09.156) to that
diagram while fixing both endpoint objects in their entirety.
The restricted word is internal to $\mathcal D_\omega$.
Naturality of the inclusions on every arrow, followed by
inversion when the word traverses an arrow backwards, identifies
the restricted word with the original word after localization.
The endpoint inclusions are identities. This proves fullness.

For faithfulness take two words internal to $\mathcal D_\omega$
that become equal in the ambient all-arrow localization. Its
presentation uses identities, composition of actual arrows,
and cancellations of an arrow with its formal inverse. The
given equality therefore has a finite sequence of these
relations as a witness. Take the finite diagram of all objects
and actual arrows occurring in this witness, and apply (S09.156).
Fix in their entirety every countable object occurring in either
of the two original words. Each composition relation restricts
to the same equality of actual maps; identity and inverse
cancellation relations remain valid. Hence the restricted
finite witness proves the same equality inside the smaller
localization, with both original words unchanged. This proves
faithfulness and the two claimed equivalences. Composing the
equivalence for $\mathcal C^{[d]}_\omega$ with A-S09-24 proves
the terminal-localization reduction.

For the common-source reduction, first use A-S09-27 to replace
an arbitrary pair $A,B$ by $E_A,E_B\in\mathcal E^{[d]}$ with
the same first skeleton $P(N)$ and actual maps to $A,B$.
Take countable subobjects $E'_A\hookrightarrow E_A$ and
$E'_B\hookrightarrow E_B$ by (S09.156). Their entire first
skeletons are retained, so they still have that same $P(N)$.
A source internal to $\mathcal E^{[d]}_\omega$ for this pair,
composed with the displayed inclusions and replacement maps,
gives an actual source for $A,B$. In the reverse direction,
begin with a pair in $\mathcal E^{[d]}_\omega$ and an arbitrary
admissible common source. Apply A-S09-27 to replace that source
by a source in $\mathcal E^{[d]}$, and then take its countable
subobject by (S09.156). Composing its actual arrows gives an
internal countable common source. This proves both directions,
including the restricted test on equal target cutoffs.

Finally we give an object witnessing both proper inclusions.
Choose integers $b\geq d+1$ and $N>b$ with $N\geq2$ large
enough for the standard-theory replacement in A-S09-27.
Let $E_0$ be that replacement of $T_NG$, using the edges
$x_w\mapsto\gamma_w(x)$. It is a countable object of
$\mathcal E^{[d]}$ with support and flat terms: all standard
terms and free disk algebras are countable, and every matching
diagram and pullback in a fixed simplex is a finite limit of
countable groups. Its weight-zero algebra is the constant
$\mathbb Z$, so weight-zero projection is an actual augmentation
$\alpha:E_0\to\mathbb Z$.

Take an uncountable set $I$ and put
$V=\bigoplus_{i\in I}\mathbb Z$. Let $M=K(V,2)$ be placed
in weight $b$ and form degree $d+1$, with zero vertical
differential. Its first two simplex terms vanish. Form the
actual square-zero extension
\[
 E_*=E_0\ltimes M,\qquad
 (a,m)(a',m')=(aa',\alpha(a)m'+\alpha(a')m),\qquad M^2=0.
 \tag{S09.160}
\]
Use the old differential on $E_0$ and zero on $M$. On even
elements of $M$ put $\gamma_1(m)=m$ and
$\gamma_r(m)=0$ for $r\geq2$; on mixed even elements use the
PD addition formula. Positive-weight coefficients act as zero
on $M$, so the mixed positive-weight PD terms vanish. These
formulas give the strict augmented PD algebra, with split
inclusion and projection to $E_0$, exactly as verified in
(S09.124). The form degree $d+1$ is positive and has zero
differential, so the same verification includes odd form degree
using the required alternating square-zero product.

The first skeleton is unchanged. The additional cohomology is
only in form degree $d+1$, so comparison through $d$ remains
valid in every weight. The extra row occurs only in weight
$b<N$, leaving all rows through $d+2$ exact from $N$.
All eventual bounds may be increased beyond $b$ for the other
requirements. The support condition follows from $b\geq d+1$,
and all terms are free or torsionfree abelian groups. Hence
$E_*\in\mathcal E^{[d]}\subset\mathcal C^{[d]}$ in every
stated variant. Its simplex-two group in weight $b$, form
$d+1$, contains $V$ as an actual direct summand. It is
uncountable and cannot be strictly isomorphic to a countable
object. This proves both properness assertions.
\end{proof}
